\phantomsection
\subsection*{Aroma e Sabor}
\addcontentsline{toc}{subsection}{Aroma e Sabor}

\begin{itemize}

\item As \textbf{variedades de mel} utilizadas podem apresentar um caráter de sutil a intenso, próprio de cada variedade (floral, frutado, condimentado, terroso etc.), perceptível no aroma e no sabor do hidromel. O hidromel não deve ter sabor de um coquetel feito com mel cru. A intensidade da percepção do mel geralmente aumenta com a força alcoólica e o dulçor.
\item No mundo do vinho, os aromas normalmente são classificados de acordo com sua origem: aqueles provenientes dos ingredientes, aqueles formados durante a fermentação e a maturação e aqueles desenvolvidos durante o envelhecimento. Juízes de hidromel podem se beneficiar de uma abordagem semelhante, compreendendo como o processo de produção transforma os ingredientes originais em um produto finalizado. O buquê resultante do processo geralmente se torna mais perceptível em hidroméis mais leves, secos ou delicados.
\item Um perfil de fermentação limpo é desejável na maioria dos estilos, mas isso não significa ausência total do \textbf{caráter da levedura}. A levedura pode acrescentar notas de ésteres ou um leve e agradável frescor; essas características não são defeitos. A maturação com a presença de leveduras pode contribuir com leves notas de nozes, tostado ou pão.
\item Notas de enxofre com aroma desagradável, dominantes ou excessivas são um defeito e normalmente indicam uma fermentação inadequada. A oxidação é indesejada na maioria dos estilos (exceto naqueles em que o envelhecimento prolongado é típico) e frequentemente se manifesta como perda de intensidade e vivacidade, aromas e sabores apagados ou uma leve lembrança de nozes semelhante à de Jerez. No entanto, um caráter agradável de envelhecimento é desejável em muitos hidroméis de maior força alcoólica.
\item O \textbf{dulçor} residual percebido no hidromel auxilia na sua classificação para fins de julgamento. A quantidade de açúcar residual medida ou a densidade final não são os únicos fatores determinantes — taninos, acidez, densidade inicial, álcool e outros ingredientes podem alterar essa percepção.

\begin{itemize}

\item \textbf{Seco (Dry)}: FG 0,990 - 1,010. Sem percepção de dulçor, mas não precisa apresentar uma sensação extremamente seca.
\item \textbf{Médio (Semi-sweet)}: FG 1,010 - 1,025. A percepção é doce, mas ainda equilibrada e refrescante. Também denominado dulçor médio.
\item \textbf{Doce (Sweet)}: FG 1,025 - 1,050. De claramente doce até níveis semelhantes aos de vinhos de sobremesa. Não deve ser como xarope, enjoativo ou apresentar sabor de mel não fermentado. Alguns estilos regionais e hidroméis de sobremesa podem ultrapassar esse nível.

\end{itemize}

\item Durante o julgamento, organize as amostras em ordem crescente de dulçor e de impacto no paladar (incluindo a picância). Tenha em mente que o dulçor pode mascarar defeitos, portanto, esteja mais atento a eles nos hidroméis mais doces. Da mesma forma, não penalize excessivamente hidroméis secos por pequenos defeitos que podem estar perceptíveis apenas devido à ausência de dulçor.
\item A \textbf{acidez} é um elemento essencial no equilíbrio do dulçor, estimula a salivação e proporciona uma impressão limpa, viva, fresca, vibrante, suculenta e refrescante, sem causar sensação de repuxamento na boca. Notas de solvente ou vinagre são defeitos indesejáveis. O dulçor sem equilíbrio de acidez resulta em uma bebida sem vivacidade e desinteressante, descrita como \textbf{enjoativa} \textit{(cloying)}.
\item Os \textbf{taninos} fornecem adstringência, corpo e, algumas vezes, amargor, contribuindo para o equilíbrio, a estrutura e a facilidade de consumo. Os taninos podem fazer com que um hidromel pareça mais seco do que seu nível de açúcar residual poderia sugerir. A percepção dos taninos pode variar de macia, suave e aveludada até áspera, abrasiva e marcante.
\item O \textbf{álcool} pode ser percebido nos exemplares de maior força alcoólica como uma leve sensação condimentada, mas não deve ser pungente, semelhante a solvente, agressivo ou provocar uma sensação intensa de queimação.
\item Taninos e acidez atuam em conjunto para equilibrar o dulçor, o caráter do mel e o álcool.

\end{itemize}